\chapter*{Введение}
\label{sec:afterwords}
\addcontentsline{toc}{chapter}{Введение}

Изучение исторических документов играет важную роль в понимании культурной,
социальной и политической истории страны. В последние годы возрос интерес к разработке
методов автоматического анализа исторических документов с использованием техник
машинного обучения и компьютерного зрения. В НИЯУ МИФИ открылась лаборатория
цифровой лингвистики в сотрудничестве со специалистами Института русского языка им.
В.В. Виноградова РАН в целях разработки новых методов цифрового анализа древнерусских
рукописей. Основная цель лаборатории — создание и поддержка корпуса старославянского
языка. Корпус языка используется для изучения различных аспектов языка, таких как
грамматика, лексика, синтаксис, стилистика, дискурс и т. д. Для лингвистов корпус является
основным инструментом поиска примеров того или иного феномена в языке. В целом, наша
работа направлена на создание новых инструментов и методов для автоматического анализа
исторических старославянский текстов.

Задача сегментации и классификации крюков в старославянских текстах является уникальной и сложной, 
поскольку первоочередная важность у проектов, подобных корпусу, заключается в классификации текстов. 
В результате, зачастую отсутствуют готовые решения для таких специфических задач. Специфичность данной 
задачи также накладывает ключевое ограничение — отсутствие датасетов или обширных выборок для её решения. 
Таким образом, для успешного выполнения задачи классификации и сегментации крюков необходимо, 
в первую очередь, собрать обширную выборку крюков, провести её корректировку с участием лингвистов 
и создать первую версию нейронной модели для сегментации крюков.

Дополнительной сложностью является разность крюков в старославянских текстах по времени: 
крюки претерпели значительные изменения, и крюки середины XVII века значительно отличаются от крюков XII - XVI веков. 
Более того, из-за изменения стандартов крюков многие рукописи переписывались поверх прежних 
текстов, что делает задачу ещё сложнее.

Наша работа посвящена созданию базовых средств для классификации крюков 
середины XVII века и последующему запуску нейронной модели для дальнейшей сегментации.

%%% Local Variables:
%%% TeX-engine: xetex
%%% eval: (setq-local TeX-master (concat "../" (seq-find (-cut string-match ".*-3-pz\.tex$" <>) (directory-files ".."))))
%%% End:
